% ======================================================================
% doc_body.tex -- the main document content.
%
% This template is for UNCG's PHY 419 (Advanced Lab) course,
% Assignment 1: Research Proposal.
%
% Adapted from the APS REVTeX 4.2 template in Overleaf on 2019-08-20.
%
% This file is part of the APS files in the REVTeX 4.2 distribution.
%     Version 4.2a of REVTeX, December 2014
%     Copyright (c) 2014 The American Physical Society.
%     See the REVTeX 4 README file for restrictions and more information.
%
% Using this requires BibTeX. Overleaf handles this automatically,
% **if** you have the right files in your project folder (*.bib),
% and you recompile both this document and the *.bib document after
% making changes to what's in *.bib.
%
% Note that the comment character "%" has been used to deactivate
% certain lines below, but keep the code in the file in case you want
% to reactivate them at some point.% to reactivate them at some point.


%
% ======================================================================
\documentclass[%
 reprint,
%superscriptaddress,
%groupedaddress,
%unsortedaddress,
%runinaddress,
%frontmatterverbose, 
%preprint,
%preprintnumbers,
%nofootinbib,
%nobibnotes,
%bibnotes,
 amsmath,amssymb,
 aps,
%pra,
%prb,
%rmp,
%prstab,
%prstper,
%floatfix,
]{revtex4-2}

\usepackage{graphicx}% Include figure files
\usepackage{dcolumn}% Align table columns on decimal point
\usepackage{bm}% bold math
\usepackage{hyperref}% enables hyperlinks in the PDF

\begin{document}

\title{Proposal Title}% Force line breaks with \\

\author{Lead Author} % Probably you, the student
  \email{lwauthor@uncg.edu}
\author{Second Author} % Probably your research mentor
 \email{Second.Author@institution.edu}
\affiliation{%
 Department of Physics \& Astronomy\\
 University of North Carolina Greensboro
}%

\date{\today}% It is always \today, today,
             %  but any date may be explicitly specified

\begin{abstract}
This document is both a brief description of the parts of a research proposal and a \LaTeX\ template to help you format one. Make a copy of the \textit{Overleaf} project, and then edit its contents to produce your proposal. The abstract, should be a terse but specific paragraph summarizing the entire document at a high level.
\end{abstract}

\maketitle
% ========================================================================

\section{Introduction \& Background}

First, set the context for the research project and provide a summary of relevant background research (by your own team and/or others) that will help the reader understand how your project fits into the larger research landscape. Cite all prior research with literature references, like this \cite{feyn54} and this \cite{Beutler1994,witten2001,Bire82,epr}, linked to a references list at the end. (Let Bib\TeX do the numbering and formatting.)

Then, introduce your particular research topic or objective, in such a way that the reader understand why it is interesting and worthwhile.

By the way, the outline I've provided in this template is not the only way to structure your proposal or break it up into sections. Addressing all these topics is important, but different people and different projects will lead to different organizational structures.

\section{Research Goal(s) \& Question(s)}

Explicitly state the goal(s) of your research project, and articulate the question(s) it will attempt to answer. If appropriate, state the research hypothesis(es) that will be tested. For a sufficiently complicated project, breaking down goals and/or questions into a hierarchy of sub-goals and/or sub-questions can be helpful.

(Sometimes, merging this section into the end of the prior sections makes more sense. It depends on how extensive the section is; use your judgment.)

\section{Research Plan}

Articulate the specific activities that will be undertaken to conduct the research project. Purpose: Convince the reader that you've thought through what has to happen in adequate detail, and that the activities planned are practical.

\subsection{Timeline}

Include a rough projection of what activities will occur when. This can be a simple bulleted list, a table, a boxes-and-arrows diagram, or whatever seems clearest and most compact. Your goal is to demonstrate to the reader that the project can reasonably be completed in the time available, and to identify milestones against which you will measure progress.

\subsection{Equipment \& Apparatus}

Identify the equipment and experimental apparatus necessary to carry the project. Standard equipment can simply be listed (and don't get too bogged down in detail); custom apparatus should be described in detail, probably with diagrams. Indicate what is already available, what must be purchased (or otherwise obtained), and what must be designed and/or fabricated. (If iteratively developing custom equipment will be a significant focus of the research project, preliminary plans are adequate.)

For examples of how to include figures or tables, as well as do various kinds of fancy formatting, see the example code in the Rev\TeX\ template available here: \url{https://tinyurl.com/y5gvss3m} (by following the link and then clicking the "Open as Template" button). The \textit{Overleaf} help system also has lots of useful information.

\section{Budget}

State total funds you'll need for the work (if any), broken down by what they're needed for. Estimates are okay. Identify where funding will come from, and be clear what is already available, what has been promised, what has been requested but not approved, and what has not yet been requested. Purpose: Convince the reader that the project's funding needs can be met. If this proposal is requesting funding support (which it's not for PHY 419), the purpose is to convince the reader that the money requested is necessary and appropriate.)

For a UNCG Undergraduate Research and Creativity Grant proposal, the budget is usually a modest paragraph or so. For a National Science Foundation proposal, it's generally a one- or two-page text document accompanied by several pages of \textit{very} detailed spreadsheets.

\section{Personnel}

This contains brief descriptions of each member of the research team --- about one paragraph each, perhaps short. Identify each person's status (undergraduate or graduate student researcher, postdoctoral research associate, faculty member, etc.) and their role in the project, and \textit{briefly} summarize their relevant qualifications or experience. Don't exaggerate or posture; just state credentials. Your purpose here is to convince the reader that your team has adequate expertise for the project.

\section{References}

At the end of the document, include the list of references cited, with full bibliographic information. Citation formats vary, and proposals are usually less strict about that than journals. It's hard to go wrong with the format used by a major journal in the discipline.

BTW, don't try to format citations by hand. You'll get stuff wrong and probably go insane in the process. Learn to use Bib\TeX. If you include your reference data in a \texttt{.bib} file, \textit{Overleaf} can automagically format it for you. Many citation managers, including the free \textit{Zotero}, can export an appropriate Bib\TeX\ file for you. It's well worth your time to learn how to use one, and the UNCG Libraries has resource web pages and even help workshops to teach you \textit{Zotero}.

\bibliography{bibfile}% Produces the bibliography via BibTeX.

\end{document}

%%%%%% End of file doc_body.tex %%%%%%

% ========================================================================
